\section*{Chapitre \ref{ch:qm:stochastic-control} --- Contrôle stochastique}

\begin{solution}{exo:qm:stochastic-control:1}
$0.04/(4 \times 0.04) = 25\%$.
\end{solution}

\begin{solution}{exo:qm:stochastic-control:2}
$0.02 + 0.05 - \tfrac12 \times 3 \times 0.0324 = 2.14\%$ par an, contre 3.29~\% à l'optimum.
\end{solution}

\begin{solution}{exo:qm:stochastic-control:3}
Pour $\ln w$~: $-w(-1/w^2)/(1/w) = 1$. Pour $w^{1-\gamma}/(1 - \gamma)$~: $-w(-\gamma w^{-\gamma-1})/w^{-\gamma} = \gamma$.
\end{solution}

\begin{solution}{exo:qm:stochastic-control:4}
Le rendement en \omterm{def:qm:stochastic-control:ce}{équivalent certain} $c(\pi) = r + \pi(\mu - r) - \tfrac12\gamma\pi^2\sigma^2$ est une parabole dont le maximum est atteint en $\pi^*$ et
de dérivée seconde $-\gamma\sigma^2$, donc $c(\pi^*) - c(\pi) = \tfrac12\gamma\sigma^2(\pi - \pi^*)^2 = \tfrac12 \times 3 \times 0.0324 \times 0.0856^2 = 3.6$ pb.
\end{solution}

\begin{solution}{exo:qm:stochastic-control:5}
$0.05/0.0324 = 154\%$, soit un emprunt de 54~\% de la richesse. L'aversion au risque du fonds est supérieure à celle de l'utilité logarithmique, le
rendement espéré est incertain (surpondérer coûte plus que sous-pondérer~; livre~2, chapitre~29), et l'effet de levier apporte
des contraintes de marge et de drawdown que le modèle ignore.
\end{solution}

\begin{solution}{exo:qm:stochastic-control:6}
Par la preuve du \cref{thm:qm:stochastic-control:merton}~: $\partial_tV = -(1 - \gamma)\rho V$ et le supremum vaut $(1 - \gamma)\rho V$,
donc l'équation est vérifiée. Pour le fonds, $\rho = 0.02 + 0.05^2/(2 \times 3 \times 0.0324) = 3.29\%$.
\end{solution}

\begin{solution}{exo:qm:stochastic-control:7}
Oui~: 0.315 à tout niveau de richesse, contre $(\mu - r)/(\gamma\sigma^2) = 0.309$ et un optimum à une période de 0.309~; le faible écart
provient de l'interpolation de la \omterm{def:qm:stochastic-control:problem}{fonction valeur} sur la grille.
\end{solution}

\begin{solution}{exo:qm:stochastic-control:8}
Dix ans de rendements à 18~\% de volatilité estiment la prime de risque avec une erreur type de $18\%/\sqrt{10} = 5.7\%$, ce qui
déplace la \omterm{def:qm:stochastic-control:fraction}{fraction de Merton} de $\pm 59$ points~: les données ne permettent pas de distinguer 0~\% de 100~\%. Réduire l'estimation vers un
a priori (chapitre~14) et dimensionner en fonction de l'incertitude plutôt que de l'estimation ponctuelle.
\end{solution}

\begin{solution}{pb:qm:stochastic-control:1}
\textbf{1.} $0.05/(3 \times 0.0324) = 51.4\%$.
\textbf{2.} $\rho = 3.29\%$.
\textbf{3.} 3.6 pb par an.
\textbf{4.} 0.02 pb~: 52~\% est presque exactement optimal.
\textbf{5.} Merton~: vendre 0.6 point, presque rien~; le règlement de marché~: racheter 8 points pour revenir à 60~\%.
\textbf{6.} 77~\% et 31~\%.
\textbf{7.} 41~\% et 62~\%.
\textbf{8.} 154~\%~: le critère de Kelly, fraction optimale en croissance, avec emprunt.
\textbf{9.} 0.52 à tout niveau de richesse et à toute date~; l'optimum à une période est 0.517 (en continu~: 0.514).
\textbf{10.} L'\omterm{def:qm:stochastic-control:utility}{utilité CRRA} est homothétique~: multiplier la richesse multiplie l'utilité, si bien que la meilleure fraction est la même à chaque
niveau de richesse.
\textbf{11.} 3.28~\%, 3.31~\% et 3.24~\% (théorie de Merton~: 3.29~\%).
\textbf{12.} 50~\%, 68~\% et 81~\%.
\textbf{13.} 3.3 pb par an.
\textbf{14.} 5.7~\%~; $\pi^*$ se déplace de $0.057/(3 \times 0.0324) = 59$ points dans les deux sens, soit d'environ $-7\%$ à $110\%$.
\textbf{15.} L'incertitude sur $\mu$, de loin~: le comité débat de quelques points de base d'\omterm{def:qm:stochastic-control:ce}{équivalent certain}.
\textbf{16.} Une bande autour de la cible à l'intérieur de laquelle le fonds ne négocie pas (chapitre~10).
\textbf{17.} Le fonds détiendrait un portefeuille de couverture du passif, plus une position de type Merton dont la taille dépend de son
excédent de financement~; la fraction dépendrait de la richesse.
\textbf{18.} Rééquilibrer mécaniquement vers une cible~; consacrer le débat aux paramètres (rendement espéré, aversion au risque) et
aux coûts, et non au timing.
\textbf{19.} Résultat à retenir~: pour \emph{le comité de rééquilibrage}, la \omterm{def:qm:stochastic-control:fraction}{fraction de Merton} vaut 51.4~\%~; conserver la cible de 60~\% du règlement de
marché coûte 3.6 pb par an de rendement en \omterm{def:qm:stochastic-control:ce}{équivalent certain}, ne jamais rééquilibrer environ 3.3 pb, et une erreur type sur
le rendement espéré déplace la fraction de 59 points.
\textbf{20.} Rééquilibrer vers la fraction constante, puisque rien dans l'optimum n'a changé avec le cours.
\end{solution}

\begin{solution}{iq:qm:stochastic-control:1}
$\pi^* = (\mu - r)/(\gamma\sigma^2)$~: le rendement espéré en excès (plus de rémunération, plus d'actions), l'aversion au risque (plus d'aversion,
moins), et la variance (plus de risque, moins)~: un compromis moyenne-variance par unité de richesse.
\iqlookfor{la formule et sa statique comparative.}
\end{solution}

\begin{solution}{iq:qm:stochastic-control:2}
$V(t, x) = \sup_u\E[\int_t^{t+h}f\,ds + V(t + h, X_{t+h})]$~; développer $V(t + h, X_{t+h})$ par Itô, prendre les espérances, soustraire $V(t,
x)$, diviser par $h$ et faire tendre $h \to 0$ : $\partial_tV + \sup_u\{\mathcal L^uV + f\} = 0$.
\iqlookfor{Itô dans le principe de la \omterm{def:qm:stochastic-control:dp}{programmation dynamique}.}
\end{solution}

\begin{solution}{iq:qm:stochastic-control:3}
La fraction optimale est constante~; une baisse fait passer la part en actions sous ce niveau, si bien que la rétablir suppose d'acheter. Rien dans
l'ensemble des opportunités du modèle n'a changé avec le cours~: les rendements sont indépendants et identiquement distribués.
\iqlookfor{fraction constante et rendements i.i.d.}
\end{solution}

\begin{solution}{iq:qm:stochastic-control:4}
Le critère de Kelly maximise l'espérance du logarithme de la richesse~: Merton avec $\gamma = 1$, soit $(\mu - r)/\sigma^2$. Les investisseurs plus averses au risque détiennent une
fraction $1/\gamma$ de celle-ci~: le Kelly fractionnaire est Merton avec $\gamma > 1$.
\iqlookfor{utilité logarithmique et Kelly fractionnaire.}
\end{solution}

\begin{solution}{iq:qm:stochastic-control:5}
Cela peut venir de ce que l'objectif est plat à cet endroit (le choix importe alors peu), de ce que la grille de contrôle est trop grossière, de ce que la
quadrature est trop rudimentaire, ou de ce que l'interpolation de la \omterm{def:qm:stochastic-control:problem}{fonction valeur} est trop grossière~; raffiner chacun de ces points et vérifier si la politique
converge.
\iqlookfor{platitude contre erreur numérique, et test de convergence.}
\end{solution}

\begin{solution}{iq:qm:stochastic-control:6}
Lorsque le contrôle optimal bascule entre des valeurs extrêmes, ou à la frontière d'arrêt, ou avec une diffusion dégénérée~: la
\omterm{def:qm:stochastic-control:problem}{fonction valeur} présente des angles. C'est alors l'unique \omterm{def:qm:stochastic-control:viscosity}{solution de viscosité}~: en chaque point, les fonctions régulières qui la touchent
par-dessus ou par-dessous satisfont les inégalités de l'équation, ce qui est aussi la notion sous laquelle les schémas sur grille convergent.
\iqlookfor{fonctions test touchant par-dessus et par-dessous.}
\end{solution}
